\chapter{Dokumentation Struktur und Beziehungen im CLD}
\label{ch:dokumentation_struktur}

Dieses Kapitel beschreibt die Struktur der Beziehungen im CLD 

% --- Zweite Ebene: Sektion ---
\section*{Definition der Subsysteme}
\label{sec:subsysteme}

% --- Dritte Ebene: Subsektion für die Beziehung ---
\subsection{Economic Growth $\rightarrow$ Income}
\label{subsec:growth_income}

Wirtschaftliches Wachstum erhöht das durchschnittliche Einkommensniveau in einer Volkswirtschaft. 

\begin{description}
    \item[Wirkungslogik:] Die Beziehung wird zunächst als annähernd linear modelliert. Keine expliziten Schwellenwerte oder Sättigungen werden modelliert.
    \item[Vorzeichen:] $+$ (positiv)
    \item[Zeit:] Die Wirkung erfolgt mit kurz- bis mittelfristiger Verzögerung, da Einkommensveränderungen über Arbeitsmarkt- und Produktivitätsdynamiken entstehen.
    \item[Annahmen:] \vspace{-0.5em} % Verhindert zu großen Abstand zur Liste
    \begin{enumerate}
        \item Wirtschaftswachstum wirkt positiv auf Einkommen. 
        \item Verteilungsfragen werden explizit nicht modelliert.
    \end{enumerate}
    \item[Evidenzbasis:] Theoretisch begründet (makroökonomischer Standardmechanismus).
    \item[Unsicherheitsgrad:] Niedrig
\end{description}

